



\usepackage[left=1in,right=1in]{geometry}
\usepackage{float}
\usepackage{enumitem}
\usepackage{fancyhdr}
\usepackage{pifont}
\usepackage{commath}
\usepackage{mathtools}
\usepackage[scr]{rsfso}
\usepackage[makeroom]{cancel}
\usepackage{fancybox,framed}
\usepackage{stackengine}
\usepackage{amssymb,amsmath}
\usepackage{amsthm}
%\usepackage{unicode-math}
\usepackage{esint}
\usepackage[overload]{empheq}
\usepackage{cases} 
\usepackage{accents} 
\usepackage[symbol]{footmisc}
\usepackage{cancel}
\usepackage{wrapfig}
\usepackage{graphicx}
\usepackage{subfigure}
\usepackage{xcolor}
\graphicspath{ {figures/} }
\usepackage[titletoc]{appendix}

\usepackage{pdfpages}

\newcommand{\q}[1]{``#1''}


\renewcommand{\thefootnote}{\fnsymbol{footnote}}






\theoremstyle{definition}
\newtheorem{definition}{Definition}[section]
\newtheorem*{example}{Example}
\newtheorem*{exercise}{Exercise}
\newtheorem*{solution}{Solution}
\newtheorem*{notation}{Notation}
\newtheorem*{note}{Note}
\newtheorem*{remark}{Remark}

%\newtheorem*{claimproof}{Proof of claim \theclaim}
%\newtheorem*{theoremproof}{Proof of theorem \thetheorem}






\theoremstyle{plain}
\newtheorem{theorem}{Theorem}[section]
\newtheorem{lemma}{Lemma}[section]
\newtheorem{claim}{Claim}[section]
\newtheorem{axiom}{Axiom}[section]



%\usepackage[thmmarks]{ntheorem}


%\theorembodyfont{\normalfont}
%\newtheorem{definition}{Definition}[section]
%\theoremstyle{plain}
%\theorembodyfont{\normalfont}
%\newtheorem{example}{Example}
%\theoremstyle{nonumberplain}
%\newtheorem{notation}{Notation}
%
%\theorembodyfont{\itshape}
%\theoremstyle{plain}
%\newtheorem{theorem}{Theorem}
%\newtheorem{lemma}[theorem]{Lemma}
%
%\newtheorem{claim}{Claim}[theorem]
%\theoremstyle{nonumberplain}
%\theoremsymbol{\ensuremath{\Box}}
%\newtheorem{proof}{Proof}
%\theoremsymbol{\ensuremath{\blacksquare}} 
%\newtheorem{claimproof}{Proof of claim \theclaim}

% argumented matrix
\makeatletter
\renewcommand*\env@matrix[1][*\c@MaxMatrixCols c]{%
  \hskip -\arraycolsep
  \let\@ifnextchar\new@ifnextchar
  \array{#1}}
\makeatother
%

% volume
\makeatletter
\DeclareRobustCommand{\volume}{\text{\volumedash}V}
\newcommand{\volumedash}{%
  \makebox[0pt][l]{%
    \ooalign{\hfil\hphantom{$\m@th V$}\hfil\cr\kern0.08em--\hfil\cr}%
  }%
}
%
\makeatother



\title{%
Tensor Analysis with Applications in Engineering \\
  \large ME 5121 @ NCKU ME\\
  Lecture notes}

\author{Prof. Yi-Chun Wang}
\date{February 1, 2021}







\usepackage{hyperref}
\hypersetup{hidelinks,
	colorlinks=true,
	allcolors=black,
	pdfstartview=Fit,
	breaklinks=true
}

%------------Command----------------------------------------

\newcommand{\vc}[3]{\overset{#2}{\underset{#3}{#1}}}

%--------------tensor form---------------------
\stackMath
\newcommand\tenq[2][1]{%
 \def\useanchorwidth{T}%
  \ifnum#1>1%
    \stackunder[0pt]{\tenq[\numexpr#1-1\relax]{#2}}{\scriptscriptstyle\sim}%
  \else%
    \stackunder[1pt]{#2}{\scriptscriptstyle\sim}%
  \fi%
}

%-----------------------------------------

\DeclareMathOperator\erf{erf}
\DeclareMathOperator\erfc{erfc}

\DeclareMathOperator\Arg{Arg}
\DeclareMathOperator\Log{Log}

\DeclareMathOperator\arcsec{arcsec}
\DeclareMathOperator\arccot{arccot}
\DeclareMathOperator\arccsc{arccsc}
\DeclareMathOperator\tr{tr}

\def\Res{\mathop{\text{Res}}\limits}